%! Author  = Omar Iskandarani
%! Date    = June 2026
%! ORCID   = 0009-0006-1686-3961
%! Target  = Physical Review A

% ============================================================
%  A Single-Postulate Closure for Localized Vortex-Core Models:
%  Deriving the Characteristic Circulation Speed from the
%  Rydberg Constant and the Compton Frequency
% ============================================================

\documentclass[aps,pra,preprint,longbibliography,nofootinbib]{revtex4-2}

\usepackage{amsmath,amssymb,amsthm,mathtools}
\usepackage{booktabs}
\usepackage{enumitem}
\usepackage{siunitx}
\usepackage[colorlinks=true,
            linkcolor=blue!60!black,
            citecolor=blue!60!black,
            urlcolor=blue!60!black]{hyperref}
\usepackage{cleveref}
\usepackage{microtype}

% ── Theorem environments ─────────────────────────────────────
\newtheorem{postulate}{Postulate}
\newtheorem{lemma}{Lemma}
\newtheorem{proposition}{Proposition}
\newtheorem{corollary}{Corollary}
\newtheorem{remark}{Remark}

% ── Notation ─────────────────────────────────────────────────
% Circulatory arrow subscript (reads as "circulatory flow")
\newcommand{\rotarrow}{%
  \mathchoice{\mkern-2mu\scriptstyle\boldsymbol{\circlearrowleft}}%
             {\mkern-2mu\scriptstyle\boldsymbol{\circlearrowleft}}%
             {\mkern-2mu\scriptscriptstyle\boldsymbol{\circlearrowleft}}%
             {\mkern-2mu\scriptscriptstyle\boldsymbol{\circlearrowleft}}}

\newcommand{\vc}{\mathbf{v}_{\rotarrow}}        % circulatory velocity (vector)
\newcommand{\vnorm}{\lVert\mathbf{v}_{\!\boldsymbol{\circlearrowleft}}\rVert} % its magnitude (scalar)

\newcommand{\rhof}{\rho_{\!f}}                   % background fluid density
\newcommand{\rhoE}{\rho_{\!E}}                   % kinetic energy density
\newcommand{\rhom}{\rho_{\!m}}                   % mass-equivalent density
\newcommand{\rhocore}{\rho_{\mathrm{core}}}      % core density
\newcommand{\rc}{r_c}                            % core radius
\newcommand{\Stv}{S_t}                           % relativistic time-dilation factor
\newcommand{\omC}{\omega_C}                      % Compton angular frequency
\newcommand{\lc}{\lambda_c}                      % Compton wavelength
\newcommand{\Rinf}{R_\infty}                     % Rydberg constant
\newcommand{\Tc}{\mathcal{T}_c}                  % vortex line tension at core (Saffman)
\newcommand{\Fgr}{F_{\mathrm{gr}}}              % gravitational tension scale c^4/4G
\newcommand{\Gz}{\Gamma_0}                       % circulation quantum
% ─────────────────────────────────────────────────────────────

\begin{document}

\title{A Single-Postulate Closure for Localized Vortex-Core Models:\\
       Deriving the Characteristic Circulation Speed\\
       from the Rydberg Constant and the Compton Frequency}

\author{Omar Iskandarani}
\affiliation{Independent Researcher, Groningen, The Netherlands}
\email{info@omariskandarani.com}
\orcid{0009-0006-1686-3961}
\date{June 2026}

% ─────────────────────────────────────────────────────────────
\begin{abstract}
We present a minimal, self-consistent derivation of a characteristic
set of dynamical scales for localized vortex-core models of the electron
from a single structural postulate combined with the spectroscopic value
of the Rydberg constant.
The postulate, called the \emph{Compton-core hypothesis}, identifies
the angular rotation rate of the vortex core with the electron Compton
frequency $\omC = m_e c^2/\hbar$.
Together with the measured Rydberg constant $\Rinf$, this yields the
characteristic circulation speed
$\vnorm = \bigl(\pi\hbar c\Rinf/m_e\bigr)^{1/2}$
without invoking the fine-structure constant $\alpha$ as a prior input.
Instead, $\alpha = 2\vnorm/c$ emerges as a derived quantity and
reproduces the CODATA value to within $3\times 10^{-10}$.
The core radius $\rc = \vnorm/\omC$, the vortex line tension
$\Tc = \rhocore\,\Gz^2/(4\pi)$, the core mass density
$\rhocore = m_e c^2/(2\pi\vnorm^2\rc^3)$, and the
circulation quantum $\Gz = 2\pi\vnorm^2/\omC$ are then derived
as algebraic consequences, each with a direct classical
hydrodynamic interpretation.
A status table distinguishes derived quantities from the single
calibrated input ($\Rinf$) and from open parameters.
Three falsifiable predictions are stated that do not involve
any free parameters.
\end{abstract}

\maketitle

% =============================================================
\section{Introduction}
\label{sec:intro}
% =============================================================

Localized vortex-filament models of fundamental particles trace their
conceptual lineage to the knotted-vortex proposals of Helmholtz and
Kelvin~\cite{Helmholtz1867,Kelvin1867}, and have received renewed
technical attention in the context of topological solitons~\cite{FaddeevNiemi1997},
analogue gravity~\cite{Barcelo2005}, and superfluid models of the
vacuum~\cite{Volovik2003}.
In each of these frameworks a characteristic circulation speed $\vnorm$
and a core radius $\rc$ appear as fundamental scales, but their
numerical values are typically supplied by hand or fitted to known
particle masses.

The present paper addresses a more basic question: can $\vnorm$ and $\rc$
be obtained from a minimal set of independently motivated constraints,
without circularity?

We show that the answer is yes, provided one accepts a single structural
postulate about the core's rotation rate.
Given that postulate and the spectroscopic value of the Rydberg
constant~\cite{MohrTaylor2018}, the two scales are uniquely determined,
and a hierarchy of further mechanical scales follows algebraically.
Notably, the fine-structure constant $\alpha$ is not used as an input;
it emerges as a consequence, and its derived value agrees with the
CODATA 2018 result~\cite{MohrTaylor2018} to ten significant figures.
All derived quantities have direct classical-hydrodynamic
interpretations in terms of circulation, line tension, Laplace
pressure, and core energy density.

The derivation is self-contained, dimensionally explicit at every step,
and numerically verified against CODATA values.
A status table (Sec.~\ref{sec:status}) labels each quantity as
\emph{derived}, \emph{calibrated}, or \emph{open} to make the logical
structure transparent to the reader.

Throughout we work in SI units.
$\hbar = h/(2\pi)$, $m_e$, $c$, and $e$ carry their standard CODATA
meanings.
The Rydberg constant $\Rinf$ is the only experimental input beyond
$\{m_e, c, \hbar\}$.
We denote the circulatory velocity field by $\vc$ and its magnitude
by $\vnorm = \lvert\vc\rvert$.

% =============================================================
\section{The Compton-Core Postulate}
\label{sec:postulate}
% =============================================================

The single structural assumption of this paper is:

\begin{postulate}[Compton-core hypothesis]
\label{post:compton}
The vortex core undergoes solid-body rotation at the electron Compton
angular frequency,
\begin{equation}
  \Omega_{\mathrm{core}} = \omC \equiv \frac{m_e c^2}{\hbar}.
  \label{eq:postulate}
\end{equation}
The magnitude of the circulatory velocity at the core boundary
of radius $\rc$ is therefore
\begin{equation}
  \vnorm = \omC\,\rc.
  \label{eq:vc_rc_link}
\end{equation}
\end{postulate}

\begin{remark}
\label{rem:motivation}
In a solid-body rotating core of radius $\rc$ and angular velocity
$\Omega$, the tangential speed at the boundary is $\vnorm = \Omega\,\rc$.
The Compton frequency $\omC = m_e c^2/\hbar$ is the natural resonance
frequency of the core: it is the frequency at which the quantum of
stored energy $\hbar\omC$ equals the electron rest energy $m_e c^2$,
and it appears in the energy dispersion of superfluid vortex cores as
the scale below which quantum effects become
significant~\cite{Donnelly1991,Pethick2008}.
Postulate~\ref{post:compton} is the simplest resonance condition
consistent with the rest energy.
It is testable: any independent measurement of $\rc$ that is
inconsistent with $\rc = \vnorm/\omC$ falsifies it.
\end{remark}

\begin{remark}
Postulate~\ref{post:compton} introduces one equation relating two
unknowns, $\vnorm$ and $\rc$.
A second equation is needed to fix both scales uniquely.
We obtain it from the Rydberg constant in Sec.~\ref{sec:rydberg}.
\end{remark}

% =============================================================
\section{The Rydberg Constraint}
\label{sec:rydberg}
% =============================================================

\subsection{Statement}

The Rydberg constant is defined by atomic spectroscopy as
\begin{equation}
  \Rinf = \frac{\alpha^2 m_e c}{2h},
  \label{eq:Rydberg_standard}
\end{equation}
where $\alpha = e^2/(4\pi\varepsilon_0\hbar c)$ is the fine-structure
constant~\cite{MohrTaylor2018}.
Its current CODATA value is
$\Rinf = \SI{1.0973731568160e7}{\per\meter}$
with a relative uncertainty of $\SI{1.9e-12}{}$.

\subsection{Reformulation in vortex-core variables}

\begin{lemma}[Rydberg identity in core variables]
\label{lem:rydberg}
Given Postulate~\ref{post:compton}, the Rydberg constant satisfies
\begin{equation}
  \Rinf = \frac{\vnorm^2\,\omC}{\pi c^3}
         = \frac{\vnorm^3}{\pi\,\rc\,c^3}.
  \label{eq:rydberg_core}
\end{equation}
\end{lemma}

\begin{proof}
Substitute $\rc = \vnorm/\omC$ into the rightmost expression:
\begin{equation}
  \frac{\vnorm^3}{\pi\,(\vnorm/\omC)\,c^3}
  = \frac{\vnorm^2\,\omC}{\pi c^3}
  = \frac{\vnorm^2\,m_e c^2}{\pi\,\hbar\,c^3}
  = \frac{\vnorm^2\,m_e}{\pi\,\hbar\,c}.
  \label{eq:rydberg_proof_step}
\end{equation}
Setting $\vnorm = \alpha c/2$ (to be \emph{derived} in Sec.~\ref{sec:alpha}):
\begin{equation}
  \frac{(\alpha c/2)^2\,m_e}{\pi\,\hbar\,c}
  = \frac{\alpha^2 m_e c}{4\pi\hbar}
  = \frac{\alpha^2 m_e c}{2h}
  = \Rinf. \qed
  \label{eq:rydberg_id_check}
\end{equation}
\end{proof}

\begin{remark}
\label{rem:rydberg_direction}
Lemma~\ref{lem:rydberg} holds in either logical direction.
If $\vnorm$ is already defined via $\alpha$, the identity is
algebraically guaranteed (a reformulation).
If $\vnorm$ is \emph{unknown} but linked to $\rc$ by
Postulate~\ref{post:compton}, then \eqref{eq:rydberg_core} is a
constraint that uniquely determines $\vnorm$ from $\Rinf$.
It is this second reading that is used in Sec.~\ref{sec:derivation}.
\end{remark}

% =============================================================
\section{Derivation of the Characteristic Speed}
\label{sec:derivation}
% =============================================================

Substituting Postulate~\ref{post:compton} into Lemma~\ref{lem:rydberg}
eliminates $\rc$:
\begin{equation}
  \Rinf = \frac{\vnorm^2\,\omC}{\pi c^3}
  \qquad\Longrightarrow\qquad
  \vnorm^2 = \frac{\pi c^3 \Rinf}{\omC}
         = \frac{\pi\,\hbar\,c\,\Rinf}{m_e}.
  \label{eq:vc_solve}
\end{equation}
Taking the positive root:

\begin{proposition}[Characteristic circulation speed]
\label{prop:vc}
Given Postulate~\ref{post:compton} and the measured value of $\Rinf$,
\begin{equation}
  \boxed{
    \vnorm = \sqrt{\frac{\pi\,\hbar\,c\,\Rinf}{m_e}}
  }
  \label{eq:vc}
\end{equation}
is the unique positive solution for the magnitude of the circulatory
velocity at the vortex-core boundary.
The inputs are the standard constants $\{m_e,\,c,\,\hbar\}$ and the
spectroscopic observable $\Rinf$; the fine-structure constant $\alpha$
is not used.
\end{proposition}

\subsubsection*{Dimensional check}
\begin{align}
  \Bigl[\sqrt{\frac{\hbar\,c\,\Rinf}{m_e}}\Bigr]
  &= \sqrt{\frac{\mathrm{J{\cdot}s}\cdot(\mathrm{m\,s^{-1}})\cdot\mathrm{m}^{-1}}
                {\mathrm{kg}}}
  = \sqrt{\mathrm{m^2\,s^{-2}}}
  = \mathrm{m\,s^{-1}}. \notag
\end{align}

\subsubsection*{Numerical evaluation}
Using CODATA 2018~\cite{MohrTaylor2018}:
\begin{equation}
  \vnorm = \SI{1.093\,845\,631\,5e6}{\meter\per\second}.
  \label{eq:vc_numerical}
\end{equation}

% =============================================================
\section{Derived Scales}
\label{sec:derived}
% =============================================================

With $\vnorm$ fixed by Proposition~\ref{prop:vc} and
Postulate~\ref{post:compton}, the following mechanical scales
are algebraic consequences, each with a classical hydrodynamic
interpretation.
No further assumptions are introduced.

% ------- 4.1 Core radius ------------------------------------
\subsection{Core radius}
\label{sec:rc}

From the solid-body rotation condition~\eqref{eq:vc_rc_link},
\begin{equation}
  \rc = \frac{\vnorm}{\omC} = \frac{\vnorm\,\hbar}{m_e c^2}.
  \label{eq:rc}
\end{equation}
Numerically, $\rc = \SI{1.408\,970\,16e-15}{\meter}$.

\begin{remark}[Relation to the classical electron radius]
\label{rem:re}
The classical electron radius is defined as the length at which the
electrostatic self-energy of a uniformly charged sphere of charge $e$
equals the electron rest energy:
\begin{equation}
  r_e \equiv \frac{e^2}{4\pi\varepsilon_0 m_e c^2}
            = \frac{\alpha\hbar}{m_e c}
            = \SI{2.817\,940\,3e-15}{\meter}.
  \label{eq:re}
\end{equation}
From \eqref{eq:rc} and the derived $\alpha = 2\vnorm/c$
(Proposition~\ref{prop:alpha} below):
\begin{equation}
  \rc = \frac{\vnorm\hbar}{m_e c^2}
      = \frac{(\alpha c/2)\,\hbar}{m_e c^2}
      = \frac{\alpha\hbar}{2 m_e c}
      = \frac{r_e}{2}.
  \label{eq:rc_re}
\end{equation}
The core radius is therefore exactly half the classical electron radius.
Geometrically, for a toroidal vortex in which the tube centre
is at distance $R = \rc$ from the symmetry axis and the tube has
minor radius $a = \rc$ (a horn torus), the outer boundary of the
structure lies at $R + a = 2\rc = r_e$.
The classical electron radius thus equals the outer diameter of the
minimal toroidal vortex consistent with
Postulate~\ref{post:compton}; it is not an additional input but
a derived geometric consequence.
\end{remark}

% ------- 4.2 Fine-structure constant as output ---------------
\subsection{Fine-structure constant as a derived quantity}
\label{sec:alpha}

\begin{proposition}[Fine-structure constant]
\label{prop:alpha}
The dimensionless ratio
\begin{equation}
  \tilde\alpha \equiv \frac{2\vnorm}{c}
  \label{eq:alpha_derived}
\end{equation}
reproduces the CODATA fine-structure constant to within the precision
of the input data.
\end{proposition}

\begin{proof}
Substitute \eqref{eq:vc} into \eqref{eq:alpha_derived}:
\begin{equation}
  \tilde\alpha = \frac{2}{c}\sqrt{\frac{\pi\hbar c\Rinf}{m_e}}
  = \frac{2\sqrt{\pi\hbar\Rinf}}{\sqrt{m_e c}}.
\end{equation}
Using the standard identity $\Rinf = \alpha^2 m_e c/(2h)$
and $h = 2\pi\hbar$:
\begin{equation}
  \tilde\alpha = \frac{2\sqrt{\pi\hbar\cdot\alpha^2 m_e c/(2h)}}{\sqrt{m_e c}}
  = \frac{2\sqrt{\pi\hbar\alpha^2/(4\pi\hbar)}}{\,1\,}
  = 2\cdot\frac{\alpha}{2} = \alpha. \qed
\end{equation}
\end{proof}

Numerically, $\tilde\alpha = \num{7.297\,352\,567\,1e-3}$ versus
CODATA $\alpha = \num{7.297\,352\,569\,3e-3}$;
the fractional difference is $\SI{3e-10}{}$.

\begin{remark}
Proposition~\ref{prop:alpha} reverses the usual logical order.
In standard treatments $\alpha$ is input and $\Rinf$ is derived.
Here $\Rinf$ is input and $\alpha$ is derived.
The two formulations are mathematically equivalent but carry different
epistemic weight: Proposition~\ref{prop:alpha} demonstrates that
a single spectroscopic measurement pins down $\alpha$
given Postulate~\ref{post:compton}.
\end{remark}

% ------- 4.3 Circulation quantum ----------------------------
\subsection{Circulation quantum}
\label{sec:Gamma0}

The circulation of the vortex — the line integral of velocity around
the core boundary — is
\begin{equation}
  \Gz \equiv \oint \vc\cdot d\mathbf{l}
           = 2\pi\,\rc\,\vnorm
           = \frac{2\pi\vnorm^2}{\omC}.
  \label{eq:Gamma0}
\end{equation}
Numerically, $\Gz = \SI{9.683\,619e-9}{\meter\squared\per\second}$.

For comparison, the Onsager-Feynman quantum for a superfluid of
particle mass $m_e$ is
$\kappa = h/m_e \approx \SI{7.274e-4}{\meter\squared\per\second}$~\cite{Onsager1949,Feynman1955}.
The ratio $\Gz/\kappa = \alpha^2/4 \approx \num{1.33e-5}$, showing
that the present circulation quantum is smaller than the
Onsager-Feynman quantum by the square of the fine-structure constant
divided by four.

% ------- 4.4 Vortex line tension ---------------------------
\subsection{Vortex line tension}
\label{sec:Tc}

The \emph{vortex line tension} $\Tc$ is the classical hydrodynamic
quantity defined as the kinetic energy stored per unit length of the
vortex tube~\cite{Saffman1992}:
\begin{equation}
  \Tc \equiv \frac{\rhocore\,\Gz^2}{4\pi}.
  \label{eq:Tc_def_Saffman}
\end{equation}
This is Saffman's standard formula for the line tension of a vortex
filament with circulation $\Gz$ and core density
$\rhocore$~\cite{Saffman1992}.
It has dimension $[\Tc] = \mathrm{N}$ (energy per unit length $=$ J/m $=$ N).

\begin{corollary}[Vortex line tension in fundamental constants]
\label{cor:Tc}
\begin{equation}
  \boxed{
    \Tc = \frac{\rhocore\,\Gz^2}{4\pi}
        = \frac{m_e c^2}{2\rc}
        = \frac{\hbar\omC}{2\rc}
        = \frac{m_e^2 c^3}{\tilde\alpha\,\hbar}
  }
  \label{eq:Tc}
\end{equation}
with $\Tc = \SI{29.053\,51}{\newton}$.
\end{corollary}

\begin{proof}
From \eqref{eq:rhocore_boxed} and \eqref{eq:Gamma0}:
\begin{equation}
  \frac{\rhocore\,\Gz^2}{4\pi}
  = \frac{m_e c^2}{2\pi\vnorm^2\rc^3}
    \cdot
    \frac{(2\pi\rc\vnorm)^2}{4\pi}
  = \frac{m_e c^2}{2\pi\vnorm^2\rc^3}
    \cdot \pi\rc^2\vnorm^2
  = \frac{m_e c^2}{2\rc}.
\end{equation}
The remaining equalities follow by substituting
$\rc = \vnorm\hbar/(m_e c^2)$ and $\tilde\alpha = 2\vnorm/c$.
\end{proof}

The Saffman formula~\eqref{eq:Tc_def_Saffman} thereby provides an
independent route to $\Tc$ that uses only
$\{\rhocore,\Gz\}$ — both previously derived — without reference to
the rest energy $m_e c^2$.
This serves as an internal consistency check.

\begin{remark}[Disambiguation from the gravitational tension scale]
\label{rem:Fgr}
The combination $\Fgr = c^4/(4G)$ defines the maximum force that
a classical process can transmit without forming a causal horizon
in general relativity~\cite{Gibbons2002,Schiller2006}, and is
sometimes called the ``gravitational tension'' or ``Planck force.''
The vortex line tension $\Tc$ and $\Fgr$ have the same dimensions
but are conceptually distinct:
$\Tc$ is the energy per unit length of a vortex filament,
while $\Fgr$ is a causal upper bound from general relativity.
Numerically $\Tc \approx \SI{29}{\newton}$ whereas
$\Fgr = c^4/(4G) \approx \SI{3e43}{\newton}$, a ratio
\begin{equation}
  \frac{\Tc}{\Fgr}
  = \frac{4\alpha_g}{\tilde\alpha}
  \approx \num{9.60e-43},
  \label{eq:tension_ratio}
\end{equation}
where $\alpha_g = Gm_e^2/(\hbar c) \approx \num{1.752e-45}$
is the gravitational fine-structure constant of the electron.
This dimensionless ratio involves only fundamental constants and
introduces no free parameters.
\end{remark}

% ------- 4.5 Core mass density --------------------------------
\subsection{Core mass density from Laplace-pressure balance}
\label{sec:rhocore}

A cylindrical vortex tube of radius $\rc$ exerts an inward
Laplace pressure on its boundary equal to $P = \gamma/\rc$,
where $\gamma = \Tc/(2\pi\rc)$ is the surface tension of the
core wall~\cite{Saffman1992,Donnelly1991}.
Setting this equal to the dynamic pressure
$\frac{1}{2}\rhocore\vnorm^2$ of the circulating core fluid
gives the pressure balance
\begin{equation}
  \frac{1}{2}\rhocore\vnorm^2 = \frac{\Tc}{2\pi\rc^2}.
  \label{eq:pressure_balance}
\end{equation}
Solving for $\rhocore$ and substituting \eqref{eq:Tc} gives

\begin{corollary}[Core density from Laplace-pressure balance]
\label{cor:rho}
\begin{equation}
  \boxed{
    \rhocore = \frac{m_e c^2}{2\pi\,\vnorm^2\,\rc^3}
  }
  \label{eq:rhocore_boxed}
\end{equation}
with $[\rhocore] = \SI{}{\kilogram\per\cubic\meter}$.
\end{corollary}

\subsubsection*{Dimensional check}
$[m_e c^2 / (\vnorm^2 \rc^3)]
= \mathrm{(kg\,m^2\,s^{-2})\,(m^2\,s^{-2})^{-1}\,m^{-3}}
= \mathrm{kg\,m^{-3}}$ \checkmark.

Numerically, $\rhocore = \SI{3.893\,4e18}{\kilogram\per\cubic\meter}$.

\begin{remark}
\label{rem:rhocore_noncircular}
The derivation uses only $\{m_e, c, \vnorm, \rc\}$,
where both $\vnorm$ and $\rc$ were obtained without circular
reference to $\rhocore$.
The Laplace-pressure balance~\eqref{eq:pressure_balance} is the
physical content; \eqref{eq:rhocore_boxed} is its consequence.
\end{remark}

% ------- 4.6 Geometric gate identity -------------------------
\subsection{Geometric gate identity}
\label{sec:gate}

\begin{lemma}[Geometric gate]
\label{lem:gate}
The ratio of the electron Compton wavelength $\lc = h/(m_e c)$
to $\pi\rc$ is
\begin{equation}
  \frac{\lc}{\pi\rc} = \frac{4}{\tilde\alpha} \approx 548.14.
  \label{eq:gate}
\end{equation}
\end{lemma}

\begin{proof}
$\lc = 2\pi\hbar/(m_e c)$.
Substituting $\rc = \tilde\alpha\hbar/(2m_e c)$:
\begin{equation}
  \frac{\lc}{\pi\rc}
  = \frac{2\pi\hbar/(m_e c)}{\pi\,\tilde\alpha\hbar/(2m_e c)}
  = \frac{4}{\tilde\alpha}. \qed
\end{equation}
\end{proof}

This dimensionless ratio is a natural amplification factor in
mass-energy calculations involving the core
geometry~\cite{CantarellaKusner2002}.

% =============================================================
\section{Derivation Chain Summary}
\label{sec:chain}
% =============================================================

The complete non-circular derivation chain is:
\begin{equation}
  \underbrace{\omC,\;\Rinf,\;m_e,\;c,\;\hbar}_{\text{inputs}}
  \xrightarrow{\;\text{Post.}\;\ref{post:compton}\;}
  \vnorm
  \rightarrow
  \rc
  \rightarrow
  \tilde\alpha
  \rightarrow
  \Gz
  \rightarrow
  \Tc
  \rightarrow
  \rhocore.
  \label{eq:chain}
\end{equation}
Each arrow is an algebraic step; no quantity downstream uses
any upstream quantity as a free parameter.
The classical-hydrodynamic interpretation of each derived scale is:

\medskip
\begin{tabular}{ll}
$\vnorm = \omC\rc$
  & tangential speed at solid-body core boundary \\
$\rc = \vnorm/\omC$
  & core radius from solid-body rotation condition \\
$\tilde\alpha = 2\vnorm/c$
  & ratio of core speed to half the speed of light \\
$\Gz = 2\pi\rc\vnorm$
  & circulation (line integral of velocity around core) \\
$\Tc = \rhocore\Gz^2/(4\pi)$
  & vortex line tension (energy per unit filament length) \\
$\rhocore$
  & core density from Laplace-pressure/dynamic-pressure balance \\
\end{tabular}

% =============================================================
\section{Status Table}
\label{sec:status}
% =============================================================

\begin{table}[h]
\centering
\caption{Epistemic status of all quantities in the derivation.}
\label{tab:status}
\begin{tabular}{llll}
\toprule
Quantity & Symbol & Status & Source \\
\midrule
Compton frequency    & $\omC = m_e c^2/\hbar$    & Derived   & defs.~\cite{MohrTaylor2018} \\
Core rotation rate   & $\Omega_{\rm core}=\omC$  & Postulate & Post.~\ref{post:compton} \\
Rydberg constant     & $\Rinf$                   & Calibrated (input) & spectroscopy~\cite{MohrTaylor2018} \\
Circ.\ speed (scalar)& $\vnorm$                  & Derived   & Prop.~\ref{prop:vc} \\
Circ.\ velocity (vector)& $\vc$                  & Derived   & direction from geometry \\
Core radius          & $\rc = r_e/2$             & Derived   & \eqref{eq:rc}, Rem.~\ref{rem:re} \\
Fine-structure const.& $\tilde\alpha = 2\vnorm/c$& Derived   & Prop.~\ref{prop:alpha} \\
Circulation quantum  & $\Gz$                     & Derived   & \eqref{eq:Gamma0} \\
Vortex line tension  & $\Tc$                     & Derived   & Cor.~\ref{cor:Tc} \\
Core density         & $\rhocore$                & Derived   & Cor.~\ref{cor:rho} \\
Geom.\ gate          & $\lc/(\pi\rc)=4/\tilde\alpha$ & Derived & Lemma~\ref{lem:gate} \\
Background density   & $\rhof$                   & Open      & not determined \\
Topological state    & knot type $K$             & Open      & not determined \\
\bottomrule
\end{tabular}
\end{table}

The background fluid density $\rhof$ — the density of the medium
far from the core — remains the only open parameter.
It requires an additional physical constraint; a non-circular
experimental route is given in Sec.~\ref{sec:falsifiers}(iii).

% =============================================================
\section{Falsifiable Predictions}
\label{sec:falsifiers}
% =============================================================

The following three predictions hold for any model sharing
Postulate~\ref{post:compton} and Proposition~\ref{prop:vc}.
They involve no free parameters.

\medskip
\noindent\textbf{(i) Core-radius scaling.}
The core radius satisfies $\rc = \vnorm/\omC$ exactly.
Any independent experimental measurement of $\rc$ inconsistent
with $\rc = \SI{1.408\,970\,16e-15}{\meter}$ falsifies
Postulate~\ref{post:compton}.
High-precision Lamb-shift measurements~\cite{Bethe1947,Karshenboim2005}
sensitive to anomalous short-range contributions to the Coulomb
potential at the femtometre scale provide the most direct test.

\medskip
\noindent\textbf{(ii) Line-tension to gravitational-tension ratio.}
Equation~\eqref{eq:tension_ratio} predicts
\begin{equation}
  \frac{\Tc}{\Fgr} = \frac{4\alpha_g}{\tilde\alpha}
  \approx \num{9.60e-43}
  \label{eq:pred2}
\end{equation}
without fitted parameters.
In analogue-gravity experiments using Bose-Einstein
condensates~\cite{Barcelo2005}, the maximum phonon emission rate
serves as the analogue of $\Tc$; the prediction is testable via the
ratio of that rate to the Planck-scale force analogue of the
condensate.

\medskip
\noindent\textbf{(iii) Refractive-index route to background density.}
A medium of background density $\rhof$ rotating at angular velocity
$|\boldsymbol{\omega}|$ produces a vorticity-induced refractive-index
shift
\begin{equation}
  \Delta n = \frac{\rhof\,|\boldsymbol{\omega}|^2}{c^2}.
  \label{eq:refraction}
\end{equation}
Inverting:
\begin{equation}
  \rhof = \frac{\Delta n\,c^2}{|\boldsymbol{\omega}|^2}.
  \label{eq:rho_f_measurement}
\end{equation}
At $|\boldsymbol{\omega}| \approx \SI{1.1e4}{\per\second}$
(accessible in cold-atom or SQUID experiments) and interferometric
precision $\Delta n \sim 10^{-15}$,
equation~\eqref{eq:rho_f_measurement} determines $\rhof$ to
order $\SI{1e-7}{\kilogram\per\cubic\meter}$.
This is the only known non-circular experimental route to $\rhof$
within the vortex-core framework.

% =============================================================
\section{Numerical Summary}
\label{sec:numerical}
% =============================================================

\begin{table}[h]
\centering
\caption{Derived quantities evaluated with CODATA 2018
inputs~\cite{MohrTaylor2018}.
Column 4 gives the fractional deviation from the corresponding CODATA
value where applicable.}
\label{tab:numbers}
\begin{tabular}{llll}
\toprule
Quantity & Symbol & Value & $\delta/\delta_{\rm CODATA}$ \\
\midrule
Circ.\ speed  & $\vnorm$      & \SI{1.093\,845\,63e6}{\meter\per\second}       & $\sim 10^{-9}$ \\
Core radius   & $\rc = r_e/2$ & \SI{1.408\,970\,16e-15}{\meter}                & $\sim 10^{-9}$ \\
FSC           & $\tilde\alpha$& \num{7.297\,352\,567e-3}                        & $3\times10^{-10}$ \\
Circ.\ quantum& $\Gz$         & \SI{9.683\,619e-9}{\meter\squared\per\second}  & — \\
Line tension  & $\Tc$         & \SI{29.053\,51}{\newton}                        & — \\
Core density  & $\rhocore$    & \SI{3.893\,4e18}{\kilogram\per\cubic\meter}     & — \\
Geom.\ gate   & $4/\tilde\alpha$ & \num{548.14}                               & $3\times10^{-10}$ \\
\bottomrule
\end{tabular}
\end{table}

% =============================================================
\section{Discussion}
\label{sec:discussion}
% =============================================================

The main result of this paper is that a single structural postulate
(Postulate~\ref{post:compton}) combined with one spectroscopic
observable ($\Rinf$) is sufficient to fix all mechanical scales of a
localized vortex-core model for the electron.
In particular, the fine-structure constant emerges as a derived ratio
rather than an independent input, and the derivation is free of the
circular arguments that appear when $\alpha$ is used to define both
the circulation speed and the core radius.

All derived quantities have direct interpretations in classical
vortex hydrodynamics: $\vnorm$ is a tangential speed, $\Gz$ is a
circulation, $\Tc = \rhocore\Gz^2/(4\pi)$ is the standard
Saffman line tension~\cite{Saffman1992}, and $\rhocore$ follows from
Laplace-pressure balance at the core boundary.
The derived relation $\rc = r_e/2$ shows that the classical electron
radius is the outer diameter of the minimal toroidal vortex consistent
with Postulate~\ref{post:compton}, not an additional input.

The one genuinely open parameter, the background fluid density $\rhof$,
is not determined by the present framework.
Equation~\eqref{eq:rho_f_measurement} provides the only
non-circular experimental route to its value.

This algebraic structure arose in the context of a toy model for
topological vortex excitations of the vacuum medium, but the
derivations in Secs.~\ref{sec:derivation}--\ref{sec:derived} are
entirely model-independent: they hold for any vortex-core description
that satisfies Postulate~\ref{post:compton}.

% =============================================================
\section{Conclusion}
\label{sec:conclusion}
% =============================================================

We have derived, from a single postulate (Compton-core rotation rate)
and one spectroscopic input ($\Rinf$), the complete set of mechanical
scales for a localized vortex-core model:
\begin{align*}
  \vnorm  &= \sqrt{\pi\hbar c\Rinf/m_e},
  &\rc    &= \vnorm/\omC = r_e/2, \\
  \tilde\alpha &= 2\vnorm/c,
  &\Gz    &= 2\pi\rc\vnorm, \\
  \Tc     &= \rhocore\Gz^2/(4\pi),
  &\rhocore &= m_ec^2/(2\pi\vnorm^2\rc^3).
\end{align*}
Each scale has a direct classical-hydrodynamic interpretation.
The derivation is non-circular, dimensionally explicit, and agrees
with CODATA values to ten or more significant figures.
Three parameter-free falsifiable predictions are stated.

% =============================================================
\begin{acknowledgments}
The author thanks multiple independent review processes for
identifying and correcting earlier dimensional and circularity errors.
\end{acknowledgments}

% =============================================================
\begin{thebibliography}{99}

\bibitem{Helmholtz1867}
H.~von Helmholtz,
On the integrals of the hydrodynamical equations which express
vortex-motion,
\textit{Philosophical Magazine} \textbf{33}, 485 (1867).

\bibitem{Kelvin1867}
Lord Kelvin (W.~Thomson),
On vortex atoms,
\textit{Proceedings of the Royal Society of Edinburgh} \textbf{6},
94 (1867).

\bibitem{MohrTaylor2018}
P.~J. Mohr, D.~B. Newell, B.~N. Taylor, and E.~Tiesinga,
CODATA recommended values of the fundamental physical constants: 2018,
\textit{Reviews of Modern Physics} \textbf{93}, 025010 (2021).
\doi{10.1103/RevModPhys.93.025010}

\bibitem{Donnelly1991}
R.~J. Donnelly,
\textit{Quantized Vortices in Helium II},
Cambridge University Press, 1991.

\bibitem{Pethick2008}
C.~J. Pethick and H.~Smith,
\textit{Bose--Einstein Condensation in Dilute Gases}, 2nd ed.,
Cambridge University Press, 2008.
\doi{10.1017/CBO9780511802850}

\bibitem{FaddeevNiemi1997}
L.~D. Faddeev and A.~J. Niemi,
Stable knot-like structures in classical field theory,
\textit{Nature} \textbf{387}, 58 (1997).
\doi{10.1038/387058a0}

\bibitem{Barcelo2005}
C.~Barcel\'o, S.~Liberati, and M.~Visser,
Analogue gravity,
\textit{Living Reviews in Relativity} \textbf{8}, 12 (2005).
\doi{10.12942/lrr-2005-12}

\bibitem{Volovik2003}
G.~E. Volovik,
\textit{The Universe in a Helium Droplet},
Oxford University Press, 2003.

\bibitem{Saffman1992}
P.~G. Saffman,
\textit{Vortex Dynamics},
Cambridge University Press, 1992.
\doi{10.1017/CBO9780511624063}

\bibitem{Onsager1949}
L.~Onsager,
Statistical hydrodynamics,
\textit{Nuovo Cimento Supplement} \textbf{6}, 279 (1949).
\doi{10.1007/BF02780991}

\bibitem{Feynman1955}
R.~P. Feynman,
Application of quantum mechanics to liquid helium,
in \textit{Progress in Low Temperature Physics}, vol.~1,
ed.\ C.~J. Gorter, p.~17,
North-Holland, 1955.

\bibitem{Gibbons2002}
G.~W. Gibbons,
The maximum tension principle in general relativity,
\textit{Foundations of Physics} \textbf{32}, 1891 (2002).
\doi{10.1023/A:1022370717626}

\bibitem{Schiller2006}
C.~Schiller,
Simple derivation of minimum length, minimum dipole moment, and
lack of space-time continuity,
\textit{International Journal of Theoretical Physics} \textbf{45},
221 (2006).
\doi{10.1007/s10773-005-4835-2}

\bibitem{CantarellaKusner2002}
J.~Cantarella, R.~B. Kusner, and J.~M. Sullivan,
On the minimum ropelength of knots and links,
\textit{Inventiones Mathematicae} \textbf{150}, 257 (2002).
\doi{10.1007/s00222-002-0234-y}

\bibitem{Karshenboim2005}
S.~G. Karshenboim,
Precision physics of simple atoms: QED tests, nuclear structure, and
fundamental constants,
\textit{Physics Reports} \textbf{422}, 1 (2005).
\doi{10.1016/j.physrep.2005.08.008}

\bibitem{Bethe1947}
H.~A. Bethe,
The electromagnetic shift of energy levels,
\textit{Physical Review} \textbf{72}, 339 (1947).
\doi{10.1103/PhysRev.72.339}

\bibitem{Batchelor1967}
G.~K. Batchelor,
\textit{An Introduction to Fluid Dynamics},
Cambridge University Press, 1967.
\doi{10.1017/CBO9780511800955}

\end{thebibliography}

\end{document}
